% Options for packages loaded elsewhere
\PassOptionsToPackage{unicode}{hyperref}
\PassOptionsToPackage{hyphens}{url}
%
\documentclass[
]{article}
\usepackage{amsmath,amssymb}
\usepackage{lmodern}
\usepackage{iftex}
\ifPDFTeX
  \usepackage[T1]{fontenc}
  \usepackage[utf8]{inputenc}
  \usepackage{textcomp} % provide euro and other symbols
\else % if luatex or xetex
  \usepackage{unicode-math}
  \defaultfontfeatures{Scale=MatchLowercase}
  \defaultfontfeatures[\rmfamily]{Ligatures=TeX,Scale=1}
\fi
% Use upquote if available, for straight quotes in verbatim environments
\IfFileExists{upquote.sty}{\usepackage{upquote}}{}
\IfFileExists{microtype.sty}{% use microtype if available
  \usepackage[]{microtype}
  \UseMicrotypeSet[protrusion]{basicmath} % disable protrusion for tt fonts
}{}
\makeatletter
\@ifundefined{KOMAClassName}{% if non-KOMA class
  \IfFileExists{parskip.sty}{%
    \usepackage{parskip}
  }{% else
    \setlength{\parindent}{0pt}
    \setlength{\parskip}{6pt plus 2pt minus 1pt}}
}{% if KOMA class
  \KOMAoptions{parskip=half}}
\makeatother
\usepackage{xcolor}
\IfFileExists{xurl.sty}{\usepackage{xurl}}{} % add URL line breaks if available
\IfFileExists{bookmark.sty}{\usepackage{bookmark}}{\usepackage{hyperref}}
\hypersetup{
  pdftitle={Effect of SAE on SUA - Outcomes Upgrade (Issue \#15)},
  hidelinks,
  pdfcreator={LaTeX via pandoc}}
\urlstyle{same} % disable monospaced font for URLs
\usepackage[margin=1in]{geometry}
\usepackage{graphicx}
\makeatletter
\def\maxwidth{\ifdim\Gin@nat@width>\linewidth\linewidth\else\Gin@nat@width\fi}
\def\maxheight{\ifdim\Gin@nat@height>\textheight\textheight\else\Gin@nat@height\fi}
\makeatother
% Scale images if necessary, so that they will not overflow the page
% margins by default, and it is still possible to overwrite the defaults
% using explicit options in \includegraphics[width, height, ...]{}
\setkeys{Gin}{width=\maxwidth,height=\maxheight,keepaspectratio}
% Set default figure placement to htbp
\makeatletter
\def\fps@figure{htbp}
\makeatother
\setlength{\emergencystretch}{3em} % prevent overfull lines
\providecommand{\tightlist}{%
  \setlength{\itemsep}{0pt}\setlength{\parskip}{0pt}}
\setcounter{secnumdepth}{-\maxdimen} % remove section numbering
\usepackage{booktabs}
\ifLuaTeX
  \usepackage{selnolig}  % disable illegal ligatures
\fi

\title{Effect of SAE on SUA - Outcomes Upgrade (Issue \#15)}
\author{}
\date{\vspace{-2.5em}}

\begin{document}
\maketitle

Outcomes upgrade of issue \#15, in three parts. (A1) Net-enrollment
margin: enrollment in any higher-education institution the year after
on-time graduation, from the SIES enrollment records
(\texttt{panel\_enrollment\_2007\_2024.rds}, all institutions
2007--2024), split into the substitution margin (enrolled outside the
centralized system) and destination type (CRUCH universities, private
universities, professional institutes, technical centers), with
institutional accreditation as a quality proxy; overall and by gender.
(A2) Unconditional class ITT: the headline and gender tables
re-estimated on all 9th-grade entrants keyed to their entry class,
dropping the single-school, graduation, and on-time filters (the planned
robustness check of Appendix 7.4), plus the net-enrollment margin on the
same entrant classes. (A5) LM decomposition: the headline LM outcome
split into the taking margin (P(take exam)) and the score conditional on
taking, with the ranking-score percentile as the GPA-based sanity
outcome (issues \#8/\#10). Sample, estimator, seed and table helpers are
identical to \texttt{3\_effect\_sae\_am.Rmd}; the headline table is
replicated first as a guard.

\hypertarget{data-and-treatment}{%
\subsection{Data and treatment}\label{data-and-treatment}}

\hypertarget{higher-education-enrollment-records}{%
\subsection{Higher-education enrollment
records}\label{higher-education-enrollment-records}}

SIES enrollment records, all institutions and years 2007--2024, one
record per student x program x year. A student counts as enrolled in a
year if they hold any undergraduate (pregrado) record in it; institution
type comes from the SIES classification (CRUCH universities, private
universities, professional institutes, technical centers), and the
quality proxy is enrollment in an institution accredited in that year.
\texttt{hei\_demre} marks enrollment in a program carrying a DEMRE code
(a program admitting through the centralized system), printed below as a
cross-check of the DEMRE-based enrollment flag.

\begin{verbatim}
## panel students with any pregrado enrollment, by year:
\end{verbatim}

\begin{verbatim}
##      YEAR students  cruch  upriv     ip    cft demre_prog
##     <num>    <int>  <int>  <int>  <int>  <int>      <int>
##  1:  2007   296068 105425  89997  63590  38168      72133
##  2:  2008   405570 141518 125177  85637  54142     106055
##  3:  2009   530778 176804 166598 116974  71442     169231
##  4:  2010   655899 207020 209198 151274  89715     198279
##  5:  2011   766246 230291 249472 186425 101294     219816
##  6:  2012   840851 244790 275592 216171 105865     263460
##  7:  2013   909904 262381 287025 248453 113652     341100
##  8:  2014   955184 272901 290801 273251 119672     353639
##  9:  2015   989508 281823 294370 293361 121845     362852
## 10:  2016  1016560 293460 300646 304945 119659    1015757
## 11:  2017  1034620 303242 307651 306864 118430    1033225
## 12:  2018  1057957 316578 312594 310447 119898    1057957
## 13:  2019  1073202 346415 288285 317605 122244    1073202
## 14:  2020  1045296 356362 268860 305317 116402    1045296
## 15:  2021  1088040 376483 278400 316531 118786    1088040
## 16:  2022  1094599 373500 275819 330718 116639    1094599
## 17:  2023  1130524 377521 283090 350578 121480     660402
## 18:  2024  1159715 382326 292307 358484 128917     674433
\end{verbatim}

\hypertarget{headline-sample-and-replication-guard}{%
\subsection{Headline sample and replication
guard}\label{headline-sample-and-replication-guard}}

Same sample as the headline table: graduated on time, all four grades in
the same school, cohorts entering 9th grade 2012--2021. The six headline
outcomes replicate \texttt{3\_effect\_sae\_am.Rmd} (a check that this
file reproduces it exactly) before any new outcome is estimated.

\begin{verbatim}
## students found in any-HEI enrollment at entry + 4, by entry cohort
##  (cohorts through 2020 are covered by the SIES records):
\end{verbatim}

\begin{verbatim}
##     AGNO_9th      n enrolled_sua   any_hei   outside
##        <num>  <int>        <num>     <num>     <num>
##  1:     2012 132368    0.3013644 0.5771712 0.2805739
##  2:     2013 131771    0.3022137 0.5785188 0.2806156
##  3:     2014 125643    0.3411730 0.5850465 0.2486728
##  4:     2015 126317    0.3431130 0.5799694 0.2420102
##  5:     2016 122858    0.3312361 0.5758030 0.2485634
##  6:     2017 124962    0.3288760 0.5428370 0.2188905
##  7:     2018 124761    0.3230817 0.5310714 0.2176241
##  8:     2019 127391    0.3981521 0.5744676 0.1862847
##  9:     2020 126978    0.3916269 0.5720361 0.1883712
## 10:     2021 123067    0.4105975 0.0000000 0.0000000
\end{verbatim}

\begin{verbatim}
## 
## cross-check: DEMRE enrollment flag vs. SIES enrollment at entry + 4
##  (cohorts through 2020 ):
\end{verbatim}

\begin{verbatim}
##         enrolled_any
## enrolled      0      1
##        0 485957 268622
##        1   7043 381427
\end{verbatim}

\begin{verbatim}
## 
## DEMRE-enrolled students found in a DEMRE-coded SIES program at entry + 4: 0.982
\end{verbatim}

\begin{table}[htbp]
   \caption{Regression Results (Replication)}\label{tab: headline replication}
   \centerline{\scalebox{0.85}{\begin{tabular}{lcccccc}
      \toprule
                   & GPA & P(Take Exam) & Average LM & Applied & Assigned & Enrolled\\
                   & (1) & (2) & (3) & (4) & (5) & (6)\\
      \midrule
      Treated & -0.005 & -0.026$^{**}$ & -0.016$^{***}$ & 0.045$^{***}$ & 0.051$^{***}$ & 0.041$^{***}$\\
       & (0.018) & (0.011) & (0.004) & (0.006) & (0.005) & (0.009)\\
      \midrule
      Mean (pre-treatment) & 5.687 & 0.832 & 0.413 & 0.432 & 0.328 & 0.265\\
      Observations & 28,483 & 28,483 & 28,483 & 28,483 & 28,483 & 28,483\\
      \bottomrule
   \end{tabular}}}
\end{table}

\hypertarget{a5.-lm-decomposition}{%
\subsection{A5. LM decomposition}\label{a5.-lm-decomposition}}

The headline \texttt{Average\ LM} outcome is the within-admission-year
percentile of the average Language and Math score, with zero assigned to
students who never take the exam, so its treatment effect mixes the
taking margin with performance conditional on taking.
\texttt{P(Take\ Exam)} isolates the first margin;
\texttt{LM\ Pct.\ (takers)} the second. The ranking-score percentile is
a deterministic function of GPA relative to the school's recent history,
so its taker-conditional effect is a composition gauge for the taker
pool: it moves only if treatment changes who registers, not how they
perform (issues \#8/\#10).

\begin{table}[htbp]
   \caption{Decomposition of the LM Effect: Taking vs. Performance}\label{tab: lm decomposition}
   \centerline{\scalebox{0.85}{\begin{tabular}{lccccc}
      \toprule
                   & P(Take Exam) & Average LM & LM Pct. (takers) & Ranking Pct. & Ranking Pct. (takers)\\
                   & (1) & (2) & (3) & (4) & (5)\\
      \midrule
      Treated & -0.026$^{**}$ & -0.016$^{***}$ & -0.005 & -0.007 & -0.018$^{***}$\\
       & (0.011) & (0.004) & (0.005) & (0.008) & (0.002)\\
      \midrule
      Mean (pre-treatment) & 0.832 & 0.413 & 0.467 & 0.459 & 0.530\\
      Observations & 28,483 & 28,483 & 28,483 & 28,483 & 28,483\\
      \bottomrule
   \end{tabular}}}
\end{table}

\begin{verbatim}
## LM decomposition arithmetic: ATT(unconditional) = -0.0163;
##   taking margin  ATT(P take) x pre-mean cond. score = -0.0257 x 0.467 = -0.0120
##   performance    pre-mean P take x ATT(cond. score) = 0.832 x -0.0046 = -0.0038
##   sum of the two margins = -0.0158
\end{verbatim}

\hypertarget{a1.-net-enrollment-margin}{%
\subsection{A1. Net-enrollment margin}\label{a1.-net-enrollment-margin}}

Same headline sample and estimator, entry cohorts 2012--2020 (the last
cohort whose first admission year is covered by the SIES records).
\texttt{Enrolled\ SUA} is the headline enrollment outcome re-estimated
on these cohorts; \texttt{Any\ HEI} is enrollment in any institution in
the year after on-time graduation; \texttt{Outside\ SUA} is enrollment
in some institution without an enrollment through the centralized
system; the destination types and the accreditation proxy are
unconditional class shares, so they add up over the \texttt{Any\ HEI}
margin.

\begin{table}[htbp]
   \caption{Effect on Enrollment in Any Higher-Education Institution}\label{tab: net enrollment}
   \centerline{\scalebox{0.85}{\begin{tabular}{lcccccccc}
      \toprule
                   & Enrolled SUA & Any HEI & Outside SUA & CRUCH & Private Univ. & IP & CFT & Accredited\\
                   & (1) & (2) & (3) & (4) & (5) & (6) & (7) & (8)\\
      \midrule
      Treated & 0.028$^{*}$ & 0.035$^{***}$ & 0.009 & 0.000 & 0.035$^{***}$ & -0.003 & 0.003 & 0.034$^{***}$\\
       & (0.014) & (0.009) & (0.014) & (0.013) & (0.005) & (0.007) & (0.004) & (0.009)\\
      \midrule
      Mean (pre-treatment) & 0.265 & 0.529 & 0.268 & 0.214 & 0.112 & 0.133 & 0.070 & 0.516\\
      Observations & 25,526 & 25,526 & 25,526 & 25,526 & 25,526 & 25,526 & 25,526 & 25,526\\
      \bottomrule
   \end{tabular}}}
\end{table}

\begin{verbatim}
## two-year window (cohorts through 2019, 22,588 classes):
##   Any HEI (2y)       0.013 (0.005)   pre-mean 0.741
##   Outside SUA (2y)   0.021 (0.009)   pre-mean 0.467
\end{verbatim}

\hypertarget{by-gender}{%
\subsubsection{By gender}\label{by-gender}}

\begin{table}[htbp]
   \caption{Effect on Enrollment in Any Higher-Education Institution by Gender}\label{tab: net enrollment by gender}
   \centerline{\scalebox{0.85}{\begin{tabular}{lcccccccccc}
      \toprule
      & \multicolumn{2}{c}{Enrolled SUA}
      & \multicolumn{2}{c}{Any HEI}
      & \multicolumn{2}{c}{Outside SUA}
      & \multicolumn{2}{c}{IP or CFT}
      & \multicolumn{2}{c}{Accredited} \\
      \cmidrule(lr){2-3}\cmidrule(lr){4-5}\cmidrule(lr){6-7}\cmidrule(lr){8-9}\cmidrule(lr){10-11}
                   & Female & Male & Female & Male & Female & Male & Female & Male & Female & Male \\
                   & (1) & (2) & (3) & (4) & (5) & (6) & (7) & (8) & (9) & (10)\\
      \midrule
      Treated & 0.046$^{**}$ & 0.012 & 0.032$^{***}$ & 0.038$^{***}$ & -0.010 & 0.024 & -0.014 & 0.016$^{*}$ & 0.030$^{**}$ & 0.034$^{***}$\\
       & (0.021) & (0.010) & (0.012) & (0.013) & (0.015) & (0.016) & (0.009) & (0.009) & (0.013) & (0.012)\\
      \midrule
      Mean (pre-treatment) & 0.265 & 0.265 & 0.527 & 0.534 & 0.266 & 0.273 & 0.194 & 0.216 & 0.512 & 0.524\\
      Observations & 24,841 & 24,413 & 24,841 & 24,413 & 24,841 & 24,413 & 24,841 & 24,413 & 24,841 & 24,413\\
      \bottomrule
   \end{tabular}}}
\end{table}

\hypertarget{a2.-unconditional-entrant-classes}{%
\subsection{A2. Unconditional entrant
classes}\label{a2.-unconditional-entrant-classes}}

All 9th-grade entrants, keyed to their entry class (school x entry
year), with no single-school, graduation, or on-time requirement: the
class ITT of Appendix 7.4, on the same regular-track universe as the
headline sample. Outcomes carry zeros for students who never register
with DEMRE, so class means are over the entire entering class; GPA
averages each student's recorded year grades over the secondary years
they attended, so partial careers enter its class mean.
\texttt{Grad.\ On\ Time} makes the margin between the entrant and
headline samples visible in the same exhibit.

1,783,653 entrants; 1,517,073 graduate on time; 1,266,116 pass the
headline filters.

\begin{table}[htbp]
   \caption{Regression Results: All 9th-Grade Entrants}\label{tab: entrants}
   \centerline{\scalebox{0.85}{\begin{tabular}{lccccccc}
      \toprule
                   & Grad. On Time & GPA & P(Take Exam) & Average LM & Applied & Assigned & Enrolled\\
                   & (1) & (2) & (3) & (4) & (5) & (6) & (7)\\
      \midrule
      Treated & 0.032$^{***}$ & 0.063$^{**}$ & -0.011 & 0.000 & 0.057$^{***}$ & 0.061$^{***}$ & 0.048$^{***}$\\
       & (0.005) & (0.025) & (0.008) & (0.004) & (0.005) & (0.004) & (0.008)\\
      \midrule
      Mean (pre-treatment) & 0.819 & 5.357 & 0.758 & 0.371 & 0.383 & 0.287 & 0.231\\
      Observations & 29,360 & 29,360 & 29,360 & 29,360 & 29,360 & 29,360 & 29,360\\
      \bottomrule
   \end{tabular}}}
\end{table}

\hypertarget{by-gender-1}{%
\subsubsection{By gender}\label{by-gender-1}}

\begin{table}[htbp]
   \caption{Regression Results by Gender: All 9th-Grade Entrants}\label{tab: entrants by gender}
   \centerline{\scalebox{0.85}{\begin{tabular}{lcccccccccccc}
      \toprule
      & \multicolumn{2}{c}{GPA}
      & \multicolumn{2}{c}{P(Take Exam)}
      & \multicolumn{2}{c}{Average LM}
      & \multicolumn{2}{c}{Applied}
      & \multicolumn{2}{c}{Assigned}
      & \multicolumn{2}{c}{Enrolled} \\
      \cmidrule(lr){2-3}\cmidrule(lr){4-5}\cmidrule(lr){6-7}\cmidrule(lr){8-9}\cmidrule(lr){10-11}\cmidrule(lr){12-13}
                   & Female & Male & Female & Male & Female & Male & Female & Male & Female & Male & Female & Male \\
                   & (1) & (2) & (3) & (4) & (5) & (6) & (7) & (8) & (9) & (10) & (11) & (12)\\
      \midrule
      Treated & 0.044$^{***}$ & 0.034$^{***}$ & -0.010 & -0.027$^{***}$ & -0.001 & -0.007 & 0.077$^{***}$ & 0.022$^{**}$ & 0.072$^{***}$ & 0.030$^{***}$ & 0.064$^{***}$ & 0.020$^{**}$\\
       & (0.011) & (0.013) & (0.007) & (0.007) & (0.006) & (0.006) & (0.010) & (0.010) & (0.014) & (0.010) & (0.019) & (0.010)\\
      \midrule
      Mean (pre-treatment) & 5.445 & 5.252 & 0.788 & 0.725 & 0.372 & 0.365 & 0.417 & 0.342 & 0.296 & 0.273 & 0.233 & 0.226\\
      Observations & 28,739 & 28,276 & 28,739 & 28,276 & 28,739 & 28,276 & 28,739 & 28,276 & 28,739 & 28,276 & 28,739 & 28,276\\
      \bottomrule
   \end{tabular}}}
\end{table}

\hypertarget{net-enrollment-on-entrant-classes}{%
\subsubsection{Net enrollment on entrant
classes}\label{net-enrollment-on-entrant-classes}}

The access question on the unconditional sample: does enrollment in any
higher-education institution rise for the full entering class? Entry
cohorts 2012--2020; enrollment measured in the year after on-time
graduation (entry + 4), so late graduates enrolling later are picked up
by the two-year window printed below the table.

\begin{table}[htbp]
   \caption{Effect on Enrollment in Any Higher-Education Institution: All 9th-Grade Entrants}\label{tab: entrants net enrollment}
   \centerline{\scalebox{0.85}{\begin{tabular}{lcccc}
      \toprule
                   & Enrolled SUA & Any HEI & Outside SUA & IP or CFT\\
                   & (1) & (2) & (3) & (4)\\
      \midrule
      Treated & 0.036$^{***}$ & 0.049$^{***}$ & 0.020$^{**}$ & 0.006\\
       & (0.012) & (0.007) & (0.010) & (0.005)\\
      \midrule
      Mean (pre-treatment) & 0.231 & 0.438 & 0.223 & 0.169\\
      Observations & 26,364 & 26,364 & 26,364 & 26,364\\
      \bottomrule
   \end{tabular}}}
\end{table}

\begin{verbatim}
## two-year window (cohorts through 2019, 23,370 classes):
##   Any HEI (2y)       0.019 (0.006)   pre-mean 0.661
##   Outside SUA (2y)   0.022 (0.008)   pre-mean 0.424
\end{verbatim}

\begin{table}[htbp]
   \caption{Effect on Enrollment in Any HEI by Gender: All 9th-Grade Entrants}\label{tab: entrants net enrollment by gender}
   \centerline{\scalebox{0.85}{\begin{tabular}{lcccccccc}
      \toprule
      & \multicolumn{2}{c}{Enrolled SUA}
      & \multicolumn{2}{c}{Any HEI}
      & \multicolumn{2}{c}{Outside SUA}
      & \multicolumn{2}{c}{IP or CFT} \\
      \cmidrule(lr){2-3}\cmidrule(lr){4-5}\cmidrule(lr){6-7}\cmidrule(lr){8-9}
                   & Female & Male & Female & Male & Female & Male & Female & Male \\
                   & (1) & (2) & (3) & (4) & (5) & (6) & (7) & (8)\\
      \midrule
      Treated & 0.052$^{***}$ & 0.012 & 0.043$^{***}$ & 0.041$^{***}$ & -0.001 & 0.038$^{***}$ & -0.006 & 0.024$^{***}$\\
       & (0.019) & (0.012) & (0.014) & (0.008) & (0.012) & (0.009) & (0.007) & (0.007)\\
      \midrule
      Mean (pre-treatment) & 0.233 & 0.226 & 0.447 & 0.429 & 0.228 & 0.220 & 0.166 & 0.173\\
      Observations & 25,810 & 25,397 & 25,810 & 25,397 & 25,810 & 25,397 & 25,810 & 25,397\\
      \bottomrule
   \end{tabular}}}
\end{table}

\hypertarget{summary}{%
\subsection{Summary}\label{summary}}

\begin{verbatim}
## did version: 2.5.1 | R version 4.5.2 (2025-10-31)
\end{verbatim}

\begin{verbatim}
## -- Replication guard: committed headline table (did_by 3_effect_sae_am.Rmd) --
\end{verbatim}

\begin{verbatim}
## GPA              this run  -0.005 | draft table  -0.005 | diff  -0.0001
## Average LM       this run  -0.016 | draft table  -0.016 | diff  -0.0003
## Applied          this run   0.045 | draft table   0.045 | diff   0.0001
## Assigned         this run   0.051 | draft table   0.051 | diff  -0.0003
## Enrolled         this run   0.041 | draft table   0.041 | diff   0.0002
\end{verbatim}

\begin{verbatim}
## (draft table is the committed 5-column vintage; P(Take Exam) is the HEAD
##  6th outcome and has no committed counterpart. Bootstrap SEs under did
##  2.5.1 differ from the draft tables even though point estimates match.)
\end{verbatim}

\begin{verbatim}
## 
## -- A5. LM decomposition (headline sample) --
\end{verbatim}

\begin{verbatim}
## P(Take Exam)              -0.026 (0.011)
## Average LM                -0.016 (0.004)
## LM Pct. (takers)          -0.005 (0.005)
## Ranking Pct.              -0.007 (0.008)
## Ranking Pct. (takers)     -0.018 (0.002)
\end{verbatim}

\begin{verbatim}
## 
## -- A1. Net enrollment (headline sample, cohorts through 2020 ) --
\end{verbatim}

\begin{verbatim}
## Enrolled SUA               0.028 (0.014)
## Any HEI                    0.035 (0.009)
## Outside SUA                0.009 (0.014)
## CRUCH                      0.000 (0.013)
## Private Univ.              0.035 (0.005)
## IP                        -0.003 (0.007)
## CFT                        0.003 (0.004)
## Accredited                 0.034 (0.009)
\end{verbatim}

\begin{verbatim}
## 
##    by gender (female | male):
\end{verbatim}

\begin{verbatim}
## Enrolled SUA             female   0.046 (0.021) | male   0.012 (0.010)
## Any HEI                  female   0.032 (0.012) | male   0.038 (0.013)
## Outside SUA              female  -0.010 (0.015) | male   0.024 (0.016)
## IP or CFT                female  -0.014 (0.009) | male   0.016 (0.009)
## Accredited               female   0.030 (0.013) | male   0.034 (0.012)
\end{verbatim}

\begin{verbatim}
## 
## -- A2. Entrant classes --
\end{verbatim}

\begin{verbatim}
## Grad. On Time              0.032 (0.005)
## GPA                        0.063 (0.025)
## P(Take Exam)              -0.011 (0.008)
## Average LM                -0.000 (0.004)
## Applied                    0.057 (0.005)
## Assigned                   0.061 (0.004)
## Enrolled                   0.048 (0.008)
\end{verbatim}

\begin{verbatim}
## 
##    by gender (female | male):
\end{verbatim}

\begin{verbatim}
## GPA                      female   0.044 (0.011) | male   0.034 (0.013)
## P(Take Exam)             female  -0.010 (0.007) | male  -0.027 (0.007)
## Average LM               female  -0.001 (0.006) | male  -0.007 (0.006)
## Applied                  female   0.077 (0.010) | male   0.022 (0.010)
## Assigned                 female   0.072 (0.014) | male   0.030 (0.010)
## Enrolled                 female   0.064 (0.019) | male   0.020 (0.010)
\end{verbatim}

\begin{verbatim}
## 
##    net enrollment on entrant classes (cohorts through 2020 ):
\end{verbatim}

\begin{verbatim}
## Enrolled SUA               0.036 (0.012)
## Any HEI                    0.049 (0.007)
## Outside SUA                0.020 (0.010)
## IP or CFT                  0.006 (0.005)
\end{verbatim}

\begin{verbatim}
## 
##    by gender (female | male):
\end{verbatim}

\begin{verbatim}
## Enrolled SUA             female   0.052 (0.019) | male   0.012 (0.012)
## Any HEI                  female   0.043 (0.014) | male   0.041 (0.008)
## Outside SUA              female  -0.001 (0.012) | male   0.038 (0.009)
## IP or CFT                female  -0.006 (0.007) | male   0.024 (0.007)
\end{verbatim}

\begin{verbatim}
## 
## Guard: the replication block must match the committed headline table;
##  if it does not, the sample replication is off and the new rows are not
##  comparable to the draft.
\end{verbatim}

\end{document}
